\section*{Chapter \ref{ch:rc:model-risk-and-validation} --- Model Risk and Validation}

\begin{solution}{exo:rc:model-risk-and-validation:1}
Intended: $0.10/1.05 = 9.52\%$; divided by the sum: $0.10/2.10 = 4.76\%$, half.
\end{solution}

\begin{solution}{exo:rc:model-risk-and-validation:2}
Score $2\times3+1+1 = 8$: tier~2.
\end{solution}

\begin{solution}{exo:rc:model-risk-and-validation:3}
$\P(X\ge8)$ for a binomial with 250 trials and probability 1\%: 0.40\%. Unlikely: the yellow zone of
chapter~21, with a capital add-on.
\end{solution}

\begin{solution}{exo:rc:model-risk-and-validation:4}
With a halved VaR the exception rate rises from 1\% to perhaps a few per cent, which needs months of data
to distinguish from bad luck: at 99\% a correct model has 2.5 exceptions a year. Backtests detect gross
errors slowly; implementation tests detect them at once.
\end{solution}

\begin{solution}{exo:rc:model-risk-and-validation:5}
It tests only the implementation, not the model's assumptions: both share them. It is still the sharpest
test of the numerics (steps, boundaries, calibration), because any difference is an error of the
candidate.
\end{solution}

\begin{solution}{exo:rc:model-risk-and-validation:6}
The VaR model as a whole is a model (it applies statistical theory to data); the spreadsheet computing
relative changes is part of its implementation, even if simple arithmetic on its own is excluded. Its
control belongs to the model's validation and change management, not to a separate spreadsheet
policy only.
\end{solution}

\begin{solution}{exo:rc:model-risk-and-validation:7}
With steps of one twenty-fourth of a year, 3 of the 36 one-year points fail (against 14 with monthly
steps); the largest error is 1.40 basis points of notional, and the largest relative error, 15\%, is on
an out-of-the-money option worth little.
\end{solution}

\begin{solution}{exo:rc:model-risk-and-validation:8}
Approval certifies a review, not correctness; a model that lowers a number during a limit breach needs
more scrutiny, not less. The 2012 model was approved and wrong, and the bank's own report found the
review did not verify the implementation.
\end{solution}

\begin{solution}{pb:rc:model-risk-and-validation:1}
\textbf{1.} $(n-o)/(n+o) = \tfrac12\,(n-o)/\bigl((n+o)/2\bigr)$.
\textbf{2.} EUR~14.06 million intended, 7.01 million with the error: a ratio of 0.499.
\textbf{3.} The book is not linear in the moves (the tranche has convexity), so halving the moves does not
exactly halve the losses.
\textbf{4.} 12.0: the equity tranche's value moves about twelve times as much as the index's for a change
in the hazard rate, so it takes twelve times the notional of index protection to hedge it.
\textbf{5.} Breaches would disappear: a limit set on the old model would look half used.
\textbf{6.} Independent re-implementation of the relative-change step, or a hand calculation of one
scenario: at once.
\textbf{7.} More \omterm{def:rc:market-risk-measures:exception}{VaR exceptions} than 1\%, visible only after months.
\textbf{8.} The change arrived while limits were breached and removed the breaches, which relieved pressure
instead of prompting scrutiny of the positions.
\textbf{9.} The documentation described the intended formula; only the implementation was wrong.
\textbf{10.} Inventory of spreadsheets used in models, independent tests of their formulas, version control,
and restricted editing.
\textbf{11.} An independent validator and risk management above the desk, with a check that limits are
recalibrated.
\textbf{12.} Parallel run with the old model, reduced limits until outcomes confirm it, and review of the
implementation of every transformation of inputs.
\textbf{13.} Recalibrate them to the new model at the change, so that a model change is not a limit change.
\textbf{14.} Tier~1: it sets limits and capital for a large, complex book, with uncertain inputs.
\textbf{15.} Expertise to rebuild the calculation, independence from CIO, and standing to stop its use.
\textbf{16.} No: the trades and their size caused it; the model hid the risk and delayed the response.
\textbf{17.} That a bank must still control such arithmetic where it is part of a model's implementation;
the exclusion concerns stand-alone calculations.
\textbf{18.} Unit tests on known cases, independent re-implementation of key steps, reconciliation of
intermediate outputs, and grid benchmarking.
\textbf{19.} Named result: \emph{the averaging error}: dividing relative changes by the sum instead of the
average halves them and cuts the book's 99\% VaR from EUR~14.06 million to 7.01 million, a 50\% reduction.
\textbf{20.} The risk of acting on a model that is wrong, wrongly built or wrongly used.
\end{solution}

\begin{solution}{iq:rc:model-risk-and-validation:1}
Scope and use; conceptual soundness; data and implementation checks; benchmarking and sensitivity
analysis; \omterm{def:rc:model-risk-and-validation:bench}{outcomes analysis}; limitations; findings with severity, owners and deadlines; conditions of
use; the validator's conclusion.
\iqlookfor{structure and actionable findings.}
\end{solution}

\begin{solution}{iq:rc:model-risk-and-validation:2}
Check the model choice against the product (mean reversion, smile, number of factors); benchmark
European prices against closed forms; compare tree and regression Monte Carlo Bermudans; test
exercise boundaries and convergence; check calibration stability; compare with market prices where
available; assess the smile's effect with an alternative model.
\iqlookfor{benchmarks, convergence and model alternatives.}
\end{solution}

\begin{solution}{iq:rc:model-risk-and-validation:3}
Unit tests with known answers, limits and symmetries (parity, zero volatility, deep in or out of the
money), convergence tests, independent re-implementation, regression tests on stored outputs, and
property tests over grids.
\iqlookfor{known cases, invariants and independent implementations.}
\end{solution}

\begin{solution}{iq:rc:model-risk-and-validation:4}
Critical review by people with expertise, independence and standing. It is happening when findings
change models or their use, when validators can stop a model, and when disagreements are recorded and
escalated.
\iqlookfor{the three conditions and evidence.}
\end{solution}

\begin{solution}{iq:rc:model-risk-and-validation:5}
A new model lowered VaR during limit breaches, was approved without verifying its implementation,
and contained a spreadsheet error that halved volatility. Independent re-implementation, parallel runs,
limit recalibration at model change, and escalation of the positions behind the breaches.
\iqlookfor{the error, the governance failure and the controls.}
\end{solution}

\begin{solution}{iq:rc:model-risk-and-validation:6}
Price the product under a set of plausible alternative models and calibrations, take the dispersion
(or the distance to a conservative choice) as the model uncertainty, and reserve for it; review as
market consensus data arrive (chapter~27).
\iqlookfor{alternative models and a reserve.}
\end{solution}
